%% ==========================================================================
%%  2  Notation and Preliminaries
%%
%%  Notation follows glossary_fair_division_subsidies.md.  If a symbol is
%%  needed that the glossary does not fix, add it to the glossary first.
%% ==========================================================================
\section{Notation and Preliminaries}
\label{sec:prelim}

We study the allocation of $m$ indivisible items among $n$ agents, with
subsidies. Write $\agents = [n]$ for the set of agents and $\items$ for the set
of items, $\abs{\items} = m$. Each agent $i \in \agents$ has a valuation
$v_i : 2^{\items} \to \R$ with $v_i(\emptyset) = 0$, and we represent an
instance by the tuple $\langle \agents, \items, \set{v_i}_{i \in \agents}
\rangle$. Algorithms operate in the standard \emph{value-oracle} model: for each
$v_i$ we assume only an oracle that returns $v_i(S)$ for a queried set
$S \subseteq \items$, and running time is measured in $n$, $m$ and the number of
oracle calls.

An \emph{allocation} $\alloc = (\bundle{1}, \dots, \bundle{n})$ is an ordered
tuple of pairwise disjoint \emph{bundles}, with $\bundle{i}$ assigned to agent
$i$. The allocation is \emph{complete} if $\bigcup_{i \in \agents} \bundle{i} =
\items$ and \emph{partial} otherwise. The ordering matters throughout: much of
the theory below concerns permuting a fixed family of bundles among the agents,
and for a permutation $\sigma \in \mathbb{S}_n$ we write
$\alloc_\sigma := (A_{\sigma(1)}, \dots, A_{\sigma(n)})$ for the allocation in
which agent $i$ receives $A_{\sigma(i)}$. The utilitarian welfare of $\alloc$
is $\SW(\alloc) = \sum_{i \in \agents} v_i(\bundle{i})$.

\begin{definition}[Envy-Freeness]
\label{def:ef}
An allocation $\alloc = (\bundle{1},\dots,\bundle{n})$ is \emph{envy-free} (EF)
iff $v_i(\bundle{i}) \ge v_i(\bundle{j})$ for all agents $i, j \in \agents$.
\end{definition}

Envy-free allocations of indivisible items need not exist, so we allow a third
party to supplement the allocation with money. Agents have quasi-linear
utilities: agent $i$ receiving bundle $\bundle{i}$ and subsidy $\subsidy_i$
enjoys utility $v_i(\bundle{i}) + \subsidy_i$, with money entering linearly and
at the same exchange rate for every agent.

\begin{definition}[Envy-Free Solution]
\label{def:efsolution}
An allocation $\alloc = (\bundle{1},\dots,\bundle{n})$ and a subsidy vector
$\subsidy = (\subsidy_1,\dots,\subsidy_n) \in \R^n_{+}$ constitute an
\emph{envy-free solution} $(\alloc, \subsidy)$ iff
$v_i(\bundle{i}) + \subsidy_i \ge v_i(\bundle{j}) + \subsidy_j$ for all
$i, j \in \agents$.
\end{definition}

An allocation $\alloc$ is \emph{envy-freeable} if some $\subsidy \in \R^n_+$
makes $(\alloc,\subsidy)$ an envy-free solution; this is a property of the
allocation alone. Following~\cite{BKNS22} we write
$\Mset(\subsidy) := \set{i \in \agents \mid \subsidy_i \ge \subsidy_j \text{ for
all } j \in \agents}$ for the set of maximally subsidised agents.

We will occasionally need the standard relaxation of envy-freeness. Note that
two inequivalent forms of it appear in this literature: the goods-only form
of~\cite{HS19}, which removes one item from the envied bundle, and the two-sided
form used by~\cite{KMSTY24}, which removes one item from \emph{either} bundle.
They coincide when all items are goods and differ once chores are present, and
only the two-sided form supports the doubly monotone results. We adopt the
two-sided form.

\begin{definition}[EF1, two-sided form]
\label{def:ef1}
An allocation $\alloc$ is \emph{envy-free up to one item} (EF1) iff for all
$i, j \in \agents$ there exists $X \subseteq \bundle{i} \cup \bundle{j}$ with
$\abs{X} \le 1$ such that
$v_i(\bundle{i} \setminus X) \ge v_i(\bundle{j} \setminus X)$.
\end{definition}

% --------------------------------------------------------------------------
\subsection{Negative Dichotomous Valuations}
\label{sec:negdich}

\begin{definition}[Negative Dichotomous Valuation]
\label{def:negdich}
A valuation $v_i : 2^{\items} \to \R$ with $v_i(\emptyset) = 0$ is
\emph{negative dichotomous} iff
\[
  v_i(S) - v_i(S \cup \set{g}) \;\in\; \set{0,1}
  \qquad \text{for every } S \subseteq \items \text{ and every }
  g \in \items \setminus S,
\]
equivalently iff every marginal $\margplus{i}{S}{g} :=
v_i(S \cup \set{g}) - v_i(S)$ lies in $\set{-1,0}$.
\end{definition}

Every item is thus a chore for every agent: $v_i$ is non-increasing, so
$v_i(S) \ge v_i(T)$ whenever $S \subseteq T$. We stress that nothing further is
assumed --- the valuations need not be additive, submodular, or subadditive ---
which is the generality at which~\cite{BKNS22} works on the goods side.

It is often more convenient to work with the non-negative \emph{cost form}
\[
  \cost_i \;:=\; -\,v_i, \qquad\text{so}\qquad
  \cost_i(\emptyset) = 0, \quad
  \cost_i(S) \le \cost_i(T) \text{ for } S \subseteq T, \quad
  \cost_i(S \cup \set{g}) - \cost_i(S) \in \set{0,1}.
\]
\begin{definition}[Dichotomous cost function]
\label{def:dichcost}
A set function $\cost : 2^{\items} \to \R$ is \emph{dichotomous} if
$\cost(\emptyset) = 0$ and every marginal $\cost(S \cup \set{g}) - \cost(S)$
lies in $\set{0,1}$.
\end{definition}

Monotonicity is not a separate hypothesis: marginals in $\set{0,1}$ are
non-negative, so a dichotomous $\cost$ is automatically non-decreasing. We
nonetheless list it above because it is the property most often used directly.

\begin{remark}[Three adjectives that are routinely confused]
\label{rem:terminology}
\emph{Monotone} means only $\cost(S) \le \cost(T)$ for $S \subseteq T$, with no
bound whatever on how large an increment can be; it is far weaker than
\Cref{def:dichcost} and, on its own, supports almost none of the arguments
below. \emph{Dichotomous} adds the binary-marginal condition and is the class of
this paper. \emph{Binary additive} is the strictly smaller class
$\cost_i(S) = \abs{S \cap D_i}$ for a set $D_i$ of disliked chores, i.e.\
additive with per-item costs in $\set{0,1}$; it is the case settled
by~\cite{LMS26}.

A warning about the literature: several papers say ``binary valuations'' for
what is called dichotomous here --- \cite{BSY23} and~\cite{TWYZ23} both do ---
and reserve ``binary additive'' for the additive subclass. When a bound is
quoted, check which is meant, since the classes are not the same and the gap
between them is where this paper lives.
\end{remark}

The correspondence $v_i \leftrightarrow \cost_i$ is a bijection between negative
dichotomous valuations and dichotomous cost functions in the sense
of~\cite{BKNS22}: the class we study is exactly the pointwise negation of the
class they study. In cost form an envy-free solution is the requirement
$\cost_i(\bundle{i}) - \subsidy_i \le \cost_i(\bundle{j}) - \subsidy_j$ for all
$i,j$, read as ``my own burden net of my compensation is no worse than what I
perceive of yours''. We use whichever of $v_i$ and $\cost_i$ is clearer locally
and always state which.

The negation, however, does not carry results across, for the reasons given in
\Cref{sec:intro}. The following remark isolates the one case in which a
transformation does exist, and shows why it is uninformative.

\begin{remark}[Complementation, and why it stops at $n = 2$]
\label{rem:complement}
For a negative dichotomous $v_i$ with cost form $\cost_i$, define the
\emph{complement} $\tilde{v}_i(S) := \cost_i(\items) - \cost_i(\items \setminus
S)$. Then $\tilde{v}_i(\emptyset) = 0$ and, for $g \notin S$, writing
$T := \items \setminus (S \cup \set{g})$,
\[
  \tilde{v}_i(S \cup \set{g}) - \tilde{v}_i(S)
  \;=\; \cost_i(\items \setminus S) - \cost_i(T)
  \;=\; \cost_i(T \cup \set{g}) - \cost_i(T) \;\in\; \set{0,1},
\]
so $\tilde{v}_i$ is dichotomous in the sense of~\cite{BKNS22}. The class is
therefore closed under complementation, without any submodularity assumption.
Now let $n = 2$ and let $\alloc = (\bundle{1},\bundle{2})$ be complete, so that
$\items \setminus \bundle{2} = \bundle{1}$ and $\items \setminus \bundle{1} =
\bundle{2}$. Then $-\cost_i(\bundle{i}) = \tilde{v}_i(\bundle{j}) -
\cost_i(\items)$ and $-\cost_i(\bundle{j}) = \tilde{v}_i(\bundle{i}) -
\cost_i(\items)$, so $(\alloc, \subsidy)$ is an envy-free solution for the chores
instance $\set{v_i}$ if and only if $((\bundle{2},\bundle{1}), \subsidy)$ is an
envy-free solution for the goods instance $\set{\tilde{v}_i}$: swap the bundles
and leave each subsidy attached to its agent. Two-agent chore division is thus
exactly two-agent goods division, and~\cite{BKNS22} settles it. The step
$\items \setminus \bundle{j} = \bundle{i}$ has no analogue for $n \ge 3$, where
the complement of one bundle is a union of several others, and no reduction of
this kind is available.
\end{remark}

% --------------------------------------------------------------------------
\subsection{The Envy Graph and the Halpern--Shah Characterisation}
\label{sec:envygraph}

For an allocation $\alloc$, the \emph{envy graph} $\envygraph{\alloc}$ is the
complete weighted directed graph on vertex set $\agents$ in which the arc
$(i,j)$ carries weight
\[
  \edgew{\alloc}(i,j) \;:=\; v_i(\bundle{j}) - v_i(\bundle{i})
  \;=\; \cost_i(\bundle{i}) - \cost_i(\bundle{j}),
\]
the envy that agent $i$ has towards agent $j$. The weight of a directed path
$P$ is $\edgew{\alloc}(P) = \sum_{(i,j) \in P} \edgew{\alloc}(i,j)$, and
$\pathw{\alloc}(i)$ denotes the maximum weight of any path in
$\envygraph{\alloc}$ starting at $i$, the empty path of weight $0$ included.

The two results below are the engine of the entire subsidy line. Both are stated
by Halpern and Shah~\cite{HS19} for additive valuations, but their proofs use
only the arc weights and permutations of the bundles --- neither additivity nor
monotonicity nor any sign condition on $v_i$ enters --- and~\cite{BKNS22} invokes
them in exactly this generality for (non-additive) dichotomous valuations. They
therefore apply verbatim to negative dichotomous valuations, and we use them
without further comment.

\begin{theorem}[\cite{HS19}]
\label{thm:hs-characterisation}
For any allocation $\alloc = (\bundle{1},\dots,\bundle{n})$, the following are
equivalent.
\begin{enumerate}
  \item[(i)] $\alloc$ is envy-freeable.
  \item[(ii)] $\alloc$ maximises utilitarian welfare across all reassignments of
    its own bundles among the agents: for every permutation $\sigma$ over
    $\agents$, $\sum_{i \in \agents} v_i(\bundle{i}) \ge
    \sum_{i \in \agents} v_i(A_{\sigma(i)})$.
  \item[(iii)] The envy graph $\envygraph{\alloc}$ has no positive-weight
    directed cycle.
\end{enumerate}
\end{theorem}

\begin{theorem}[\cite{HS19}]
\label{thm:hs-minsubsidy}
For any envy-freeable allocation $\alloc$, the subsidy vector $\optsubsidy$
given by $\optsubsidy_i := \pathw{\alloc}(i)$ realises an envy-free solution
$(\alloc, \optsubsidy)$, and $\optsubsidy_i \le \subsidy_i$ for every
$i \in \agents$ and every $\subsidy$ such that $(\alloc,\subsidy)$ is an
envy-free solution. It can be computed in strongly polynomial time by running
an all-pairs longest-path computation on $\envygraph{\alloc}$.
\end{theorem}

Condition~(ii) of Theorem~\ref{thm:hs-characterisation} has the standing consequence
that an envy-freeable allocation can always be manufactured from an arbitrary
one: fix the bundles and reassign them according to a maximum-weight perfect
matching between agents and bundles. What is at issue is never the existence of
an envy-freeable allocation, only the size of the subsidy that
Theorem~\ref{thm:hs-minsubsidy} then charges for it.

Specialising to our model gives the following, which we use throughout.

\begin{observation}[Integrality]
\label{obs:integrality}
Let $\set{v_i}_{i \in \agents}$ be negative dichotomous. Then $v_i(S) \in
\Z_{\le 0}$ for every $i$ and every $S \subseteq \items$; every arc weight
$\edgew{\alloc}(i,j)$ is an integer; and for every envy-freeable allocation
$\alloc$ the minimum subsidy vector satisfies $\optsubsidy \in \Z^n_{\ge 0}$.
\end{observation}

\begin{proof}
Fix $i$ and $S = \set{g_1,\dots,g_k}$. Telescoping from $v_i(\emptyset) = 0$
along any ordering of $S$ writes $v_i(S)$ as a sum of $k$ marginals, each of
which lies in $\set{-1,0}$ by Definition~\ref{def:negdich}; hence $v_i(S) \in \Z$ and
$-k \le v_i(S) \le 0$. Consequently each arc weight
$\edgew{\alloc}(i,j) = v_i(\bundle{j}) - v_i(\bundle{i})$ is a difference of
integers, and the weight of any directed path, being a finite sum of arc
weights, is an integer. By Theorem~\ref{thm:hs-minsubsidy},
$\optsubsidy_i = \pathw{\alloc}(i)$ is the maximum such weight over paths
starting at $i$, and it is non-negative because the empty path has weight $0$.
\end{proof}

Observation~\ref{obs:integrality} is what makes the target statement --- a subsidy vector in
$\set{0,1}^n$ --- a statement about a \emph{bound} rather than about rounding:
once $\optsubsidy$ is known to be integral, showing $\optsubsidy_i \le 1$ for
every $i$ is the same as showing $\optsubsidy \in \set{0,1}^n$. The same
observation underlies the corresponding step in~\cite{BKNS22}, where the arc
weights of a dichotomous instance are likewise integers.
